\documentclass[11pt]{article}
\usepackage{amsmath,amssymb,amsthm}
\usepackage{booktabs}
\usepackage{xcolor}
\usepackage{microtype}
\usepackage[margin=1in]{geometry}
\usepackage[colorlinks=true,linkcolor=blue!55!black,citecolor=blue!55!black,urlcolor=blue!55!black]{hyperref}

\newtheorem{theorem}{Theorem}
\newtheorem{proposition}[theorem]{Proposition}
\newtheorem{lemma}[theorem]{Lemma}
\newtheorem{corollary}[theorem]{Corollary}
\theoremstyle{definition}
\newtheorem{definition}[theorem]{Definition}
\newtheorem{hypothesis}[theorem]{Hypothesis}
\newtheorem{example}[theorem]{Example}
\theoremstyle{remark}
\newtheorem{remark}[theorem]{Remark}
\newtheorem{criterion}[theorem]{Criterion}

\DeclareMathOperator{\grad}{grad}
\DeclareMathOperator{\Var}{Var}
\DeclareMathOperator{\Cov}{Cov}
\DeclareMathOperator{\KL}{KL}
\newcommand{\R}{\mathbb{R}}
\newcommand{\E}{\mathbb{E}}
\newcommand{\cS}{\mathcal{S}}
\newcommand{\cX}{\mathcal{X}}
\newcommand{\cP}{\mathcal{P}}
\newcommand{\cF}{\mathfrak{F}}
\newcommand{\cL}{\mathfrak{L}}
\newcommand{\Proj}{\Pi}

\title{\bfseries Evolutionary Equilibrium Propagation:\\[2pt]
Settling to a Stable Equilibrium as the Single Primitive of Continual Learning}
\author{Sheng Cheng\\ \small Hefei Jiuzhai Big Data Technology Co., Ltd.\\ \small \texttt{github.com/hf-jz}}
\date{October 2026}

\begin{document}
\maketitle

\begin{abstract}
\noindent
A deployed neural network today ceases to learn at deployment: its parameters are fixed at the end of
a training phase, and inference is a different kind of computation from training. Biological nervous
systems have no such boundary. We argue that this boundary characterises one design choice, the
separation of a forward computation from a backward one, rather than learning itself. We ask what a
system can achieve when its only primitive is \textbf{settling to a stable equilibrium}.

We formalise such a system as a \emph{patch net}: each unit holds local state and predicts it from its
neighbours, the network repairs disagreement by projected gradient descent on one energy, and the
settled state constitutes the prediction. Three results follow from the same operator. (i) The repair descends a
coercive real-analytic energy, so its iterates are bounded and converge to a single stationary point,
which is the precise sense in which the system \emph{always evolves toward a stable equilibrium}.
(ii) Inference and learning are the same repair restricted to different coordinate blocks, and the
learning block carries an anchored (proximal) term that makes it uniformly well posed, so that no
distinction between training and inference arises. (iii) The learning contrast is equilibrium propagation in the sense
of Scellier and Bengio, and we prove that the one-sided contrast has error $O(\beta)$ while the
\emph{centred} contrast used by current implementations has error $O(\beta^2)$, with an exactly
solvable example giving both constants in closed form.

We then lift the structure to a second time scale. Evolution acts on the \emph{meta-parameters} that
define how the system settles and learns (the wiring mask, the energy's form, the plasticity constants),
never on the weights learned within a lifetime. We prove a two-timescale averaging theorem whose effective
fitness is the post-learning (Baldwinian) fitness, a Price-type decomposition that separates innate
selection from the plasticity channel, and a free-energy dissipation identity for the population. The
construction is checked against the six formal criteria of biological evolution; the inheritance
restriction is shown to be necessary by a counterexample that is adaptive rather than merely
biological.
Every quantitative claim is reproduced by two scripts, one of which agrees with the reference
implementation to $1.1\times10^{-16}$ in float64.

\end{abstract}

\tableofcontents

\section{Introduction}

\subsection{The frozen-weight problem}
A contemporary language model is trained once, at substantial computational cost, and then used with its parameters
fixed. Online and continual variants exist, and recurrent networks can in principle continue to adapt, but
the dominant deployed regime separates a training phase from a use phase and computes gradients by
backpropagation through a stored forward pass. The separation has three consequences, none of which is easily removed. The system cannot learn from an outcome it has already observed unless a
training procedure is invoked. It cannot retain a correction across sessions unless the correction is
written back into the parameters. The correction also requires a backward pass, which no known
biological circuit performs.

Biological nervous systems do not operate in this way. A nematode with $302$ neurons and no capacity
for a backward pass navigates, learns odours, and continues to learn while it acts. This paper
therefore asks what a learning system can achieve when the forward and backward passes are replaced
by a single computation rather than optimised separately.

\subsection{The settling hypothesis}
The hypothesis we formalise is that the single primitive is \emph{settling}. A network of units, each
holding local state and predicting it from its neighbours, repairs disagreement until the whole
network is at rest. The resting state constitutes the prediction. Learning is the same repair with the
parameters made eligible, so that a single operator both answers a query and modifies the connections,
and neither a separate output head nor a backward circuit is required.

This is the design of \emph{Cadence}~\cite{cadence,manuscript}, an open research library whose
settling rule is equilibrium propagation~\cite{scellier2017}. This paper does not propose a new
library; it supplies what the design has so far lacked: a proof that settling \emph{always} reaches a
stable equilibrium under stated conditions, a proof that inference and learning are a single operator,
a second-order account of the contrast the implementation employs, and a mathematically stated
evolutionary layer with its compliance conditions.

\subsection{Contributions}
\begin{enumerate}
\item \textbf{The settling primitive is proved to converge} (\S\ref{sec:settling}). The repair energy of
a patch net with $\tanh$ predictions is coercive with or without an activity prior, because each state
enters its own prediction error and each prediction is bounded. The energy is real analytic, so the
Łojasiewicz inequality holds at every point, and a line-search descent method on such a function has
bounded iterates that converge to a single stationary point, not merely a stationary limit set.
\item \textbf{Inference and learning are one operator} (\S\ref{sec:oneop}). The learning objective
differs from the query objective by an anchored quadratic term and by which coordinate blocks are
eligible. We prove the anchored objective's parameter block is uniformly strongly convex once the
anchor weight exceeds the local curvature of the data term, which makes continual online updating
stable.
\item \textbf{The contrast is second order} (\S\ref{sec:centred}). We prove the one-sided contrast has
error $O(\beta)$ and the centred contrast, which contrasts a $+\beta$ and a $-\beta$ nudge from the
same free state, has error $O(\beta^2)$. An exactly solvable scalar instantiation gives both constants
in closed form as $\partial L/\partial w \cdot (1+\beta)^{-1}$ and
$\partial L/\partial w \cdot (1-\beta^2)^{-1}$.
\item \textbf{An evolutionary layer with an averaging theorem} (\S\ref{sec:evo}). Evolution acts on
meta-parameters. The effective fitness of a meta-parameter is the fitness after a lifetime of EP
learning, which is the Baldwinian post-learning fitness. We prove the population dynamics average to a
replicator-mutator flow with that fitness, give a Price-type decomposition, and give a free-energy
dissipation identity.
\item \textbf{Biological compliance as a theorem rather than a claim} (\S\ref{sec:bio}). Six formal criteria
of biological evolution are stated, the construction is proved to satisfy them, and the Weismann
restriction is proved necessary by a distribution of tasks on which dropping it is strictly worse.
\item \textbf{Reproducibility} (\S\ref{sec:examples}). We give four worked examples, including an exact
closed-form case and a cross-check against the reference implementation that agrees to
$1.1\times10^{-16}$.
\end{enumerate}

\subsection{Scope}
The paper proves statements about the settling dynamics and the optimisation they define. It does not
claim that settling yields intelligence, that a patch net matches a transformer at any task, or that
the evolutionary layer has been validated on hardware. \S\ref{sec:limits} states the open problems and
four falsifiable predictions. The empirical results reported here are numerical and confined to small graphs; they are evidence
about the mathematical claims and not about capability.

The results that are not new are identified as such. The learning contrast and its identification with
the gradient of the loss at infinitesimal nudge are due to Scellier and
Bengio~\cite{scellier2017} and are restated here with a proof for the patch energy. The formulation of
the same rule on a general vector field, the separation of the finite-nudge bias from the bias due to
asymmetry of the Jacobian, and the homeostatic remedy are due to Scellier et al.\ and to Laborieux and
Zenke~\cite{scellier2018,laborieux2024}, and exact gradients at finite nudge from a thermal ensemble
and from a Cauchy integral over an oscillating nudge are due to Litman and to Laborieux and
Zenke~\cite{litman2025,laborieux2022}. What is claimed here is the settling theory of
\S\ref{sec:settling} (coercivity, descent for the discrete projected rule, and convergence to a single
stationary point with a rate), the anchored parameter block and the identity of the two solves in
\S\ref{sec:oneop}, the order of the centred contrast for the patch energy in
\S\ref{sec:centred}, and the evolutionary layer of \S\ref{sec:evo} with its compliance argument in
\S\ref{sec:bio}. The convergence assumption that this literature states and leaves open is the one
supplied in \S\ref{sec:settling}.

\section{The settling primitive}
\label{sec:settling}

\subsection{The processing patch}
Let the network be a finite directed graph with $m$ sensory coordinates, $n$ patches, and an edge set
$E$ whose elements are triples $(k,s,t)$ with $k\in\{\text{input},\text{state},\text{residual}\}$.
Patch $i$ holds a live state $x_i\in[-B,B]$ and retained parameters: a weight $w_e$ per incoming edge
and a bias $b_i$. Write $s_e$ for the signal on edge $e$: a fixed sensory sample for $k=\text{input}$,
a live state $x_s$ for $k=\text{state}$, and a prediction error $e_s=x_s-p_s$ for
$k=\text{residual}$. The patch law is
\begin{equation}
p_i=\tanh\!\Big(b_i+\sum_{e=(k,s,i)\in E} w_e\,s_e\Big),
\qquad
e_i=x_i-p_i,
\label{eq:patch}
\end{equation}
and the settling energy is
\begin{equation}
E(x;\theta)=\tfrac12\sum_{i=1}^{n} e_i^{\,2}+\tfrac{\lambda}{2}\sum_{i=1}^{n} x_i^{\,2},
\qquad \theta=(w,b),\quad \lambda\ge 0 .
\label{eq:energy}
\end{equation}
Prediction and error are recomputed from state and parameters and are never independently writable,
so that they cannot conceal a disagreement. Residual edges must be acyclic, since they define derived
quantities; state edges may be recurrent, which allows a later patch to influence an earlier
one through the same energy. Sensory coordinates are boundary values, held fixed throughout a solve.

Nothing in \eqref{eq:patch} or \eqref{eq:energy} requires the connections to be symmetric, and two
different symmetries are worth separating, because only the second is used anywhere below. The first
is symmetry of the parameters. A quadratic coupling such as the Hopfield energy of the
prototype~\cite{scellier2017,hopfield} depends on the weights only through their symmetric part, so
one parameter per undirected edge is all that such an energy can see, and the source names that
requirement as the principal biological difficulty of the model. The patch energy is not a quadratic
form in the activations. Each directed edge carries its own weight and enters the energy only through
the prediction of the patch it targets, $\partial E/\partial w_e=-e_i(1-p_i^2)s_e$, so the two
directions can carry different weights and both are visible to the energy. The second symmetry is
symmetry of the Jacobian of the settling dynamics, and that one is automatic rather than assumed:
settling is a gradient flow of a scalar energy, so its Jacobian is $\nabla_x^2E$, symmetric by the
symmetry of second derivatives. The generalised formulation of equilibrium propagation on an arbitrary
vector field makes this the operative condition, its error vector agrees with the true one only when
that Jacobian is symmetric, and its asymmetry is a separate source of
bias~\cite{scellier2018,laborieux2024}.

Equation~\eqref{eq:energy} is the settling objective of the reference library~\cite{cadence}; the
implementation reproduces its value and every analytic derivative to machine precision
(Example~\ref{ex:crosscheck}). The activity prior with weight $\lambda$ is part of the model rather
than a numerical setting, and \S\ref{sec:coercive} establishes the role it plays.

\subsection{The repair operator}
Write $z$ for the vector of eligible coordinates, which are the free states for a query solve and the
free states together with the weights and biases for a learning solve. Let $\Proj$ project onto the
declared box, let $g=\nabla_z F$ be the exact gradient of the current objective. The repair operator $\mathcal{R}$ is projected gradient descent with
backtracking:
\begin{equation}
z^{(j+1)}=\Proj\big(z^{(j)}-\alpha_j\,g^{(j)}\big),
\qquad
F\big(z^{(j+1)}\big)\;\le\;F\big(z^{(j)}\big)+10^{-4}\,\big\langle g^{(j)},\,z^{(j+1)}-z^{(j)}\big\rangle ,
\label{eq:repair}
\end{equation}
with $\alpha_j$ proposed by a two-point secant (Barzilai-Borwein) estimate from the previous accepted
displacement and gradient change, halved on rejection, and with the configured initial step as the
fallback when the curvature estimate is not positive and finite. A solve qualifies when the projected
stationarity residual
\begin{equation}
\rho(z)=\big\|\,\Proj(z-g)-z\,\big\|_\infty
\label{eq:residual}
\end{equation}
is at most a declared tolerance over every eligible coordinate. Qualification is a statement about a
constrained stationary point and not about prediction error, uniqueness or accuracy: a stationary
point may retain nonzero prediction errors, and a qualified point may be inconsistent with the
underlying task.

\subsection{Descent}
In the continuous-time formulation of the source the energy identity
$\mathrm dF/\mathrm dt=-\|\nabla_xF\|^2\le0$ is immediate~\cite{scellier2017}. A discrete projected
iteration with backtracking is not a gradient flow, and the corresponding statement needs its own
argument, which is the content of the following theorem.
\begin{theorem}[Lyapunov descent]
\label{thm:descent}
Let $F$ be $C^1$ on the box, suppose every candidate step of \eqref{eq:repair} is accepted only when the
Armijo condition holds, and suppose the accepted step sizes are bounded above by $\bar\alpha<\infty$.
Then $F(z^{(j)})$ is non-increasing, the decrease is strictly negative whenever the projected residual
\eqref{eq:residual} is nonzero, and
\begin{equation}
\sum_{j}\big\|z^{(j+1)}-z^{(j)}\big\|^2\;\le\;10^{4}\,\bar\alpha\,\big(F(z^{(0)})-\inf F\big)\;<\;\infty .
\end{equation}
\end{theorem}
\begin{proof}
Write $\Delta z^{(j)}=z^{(j+1)}-z^{(j)}$ and let $I_j$ be the set of coordinates that are free after
projection, that is, those that are not at an active bound with an outward gradient. On $I_j$ the
projected step agrees with $-\alpha_j g^{(j)}$ and on its complement the displacement is zero, so
\begin{equation}
\big\langle g^{(j)},\Delta z^{(j)}\big\rangle
=-\alpha_j\sum_{c\in I_j}\big(g^{(j)}_c\big)^2 .
\end{equation}
This is strictly negative exactly when some $g_c$ with $c\in I_j$ is nonzero, which is the statement
$\rho(z^{(j)})>0$ by \eqref{eq:residual}. The Armijo inequality \eqref{eq:repair} then gives
$F(z^{(j+1)})\le F(z^{(j)})-10^{-4}\alpha_j\sum_{c\in I_j}(g^{(j)}_c)^2$. When $I_j$ is empty the step is
zero and the inequality holds with both sides equal. Otherwise
$\sum_{c\in I_j}(g^{(j)}_c)^2=\|\Delta z^{(j)}\|^2/\alpha_j^2$, so
$10^{-4}\|\Delta z^{(j)}\|^2\le\bar\alpha\big(F(z^{(j)})-F(z^{(j+1)})\big)$. Summing over $j$ and using
$F\ge0$ gives the stated bound; monotonicity and boundedness below give convergence of $F$.
\end{proof}

\subsection{Coercivity and boundedness}
\label{sec:coercive}
\begin{theorem}[Coercivity]
\label{thm:coercive}
For every $x\in\R^n$ and every $\theta$,
$F(x;\theta)\ge\tfrac12\sum_{i=1}^{n}\big((|x_i|-1)^+\big)^2$, where $t^+=\max(t,0)$. Hence
$F(x;\theta)\ge\tfrac12(\|x\|_\infty-1)^2$ whenever $\|x\|_\infty\ge1$, so $F(x;\theta)\to\infty$ as
$\|x\|\to\infty$ uniformly in $\theta$ and every sublevel set of $F$ is bounded.
\end{theorem}
\begin{proof}
Each prediction is a $\tanh$ value, so $|p_i|\le1$ and therefore $|x_i-p_i|\ge|x_i|-1$. Squaring gives
$e_i^{\,2}\ge\big((|x_i|-1)^+\big)^2$: when $|x_i|\ge1$ the right side is $(|x_i|-1)^2$ and the
inequality is the one derived above, and when $|x_i|<1$ the right side is $0$. Dropping the
nonnegative prior term and summing the remaining squared errors gives the bound, and the second
statement follows by choosing $i$ with $|x_i|=\|x\|_\infty$.
\end{proof}

Two consequences are worth stating explicitly, because the activity prior is often described as a
numerical stabiliser, which is not its principal role. First, the bound does not depend on the prior: it
follows from the architecture, because a patch always predicts its own state and the prediction is
bounded. A projected repair confines its iterates to the declared box in any case, so the content of
coercivity here is that the sublevel sets on which the descent argument runs are bounded without any
assumption on the input, which is the hypothesis Theorem~\ref{thm:convergence} uses. Second, the prior
acts by selecting among stationary points. With $\lambda=0$ the zero-error
configuration $x_i=p_i$ is stationary and attains $F=0$ when the equations are consistent, while with
$\lambda>0$ the stationary point trades prediction error against activity and typically retains a
nonzero residual. The numerical evidence is in \S\ref{ex:chain}.

\subsection{Convergence to a single equilibrium}
The hypothesis that the network \emph{always evolves toward a stable equilibrium} is stronger than
monotone descent with a stationary limit set. It requires that the trajectory converge, and the
mechanism is analyticity.

\begin{theorem}[Convergence]
\label{thm:convergence}
Assume (i) $F$ is real analytic on an open set containing the box and its sublevel sets are bounded,
(ii) every accepted step satisfies the Armijo condition \eqref{eq:repair}, and (iii) the accepted step
sizes satisfy $0<\alpha_{\min}\le\alpha_j\le\alpha_{\max}<\infty$. Then
$\sum_j\|z^{(j+1)}-z^{(j)}\|<\infty$, and the iterates converge to a single stationary point $z^\star$
with $\rho(z^\star)=0$. If $\vartheta\in[\tfrac12,1)$ is the {\L}ojasiewicz exponent of $F$ at
$z^\star$, there is a constant $C>0$ with
$\|z^{(j)}-z^\star\|\le C\,(1+j)^{-r}$ and $r=(1-\vartheta)/(2\vartheta-1)$ when $\vartheta>\tfrac12$;
convergence is exponential when $\vartheta=\tfrac12$.
\end{theorem}
\begin{proof}
A real-analytic function satisfies the {\L}ojasiewicz gradient inequality, and on a compact set that
meets only critical values the inequality holds with a uniform exponent and constant: near the critical
value $F^\star$ of the limit set there are $\vartheta\in[\tfrac12,1)$ and $c_0>0$ with
$\rho(z)\ge c_0\,|F(z)-F^\star|^{\vartheta}$~\cite{lojasiewicz}. For a feasible set defined by finitely
many analytic inequalities, which includes a box, the same inequality holds with the projected residual
$\rho$ in place of the gradient norm, by the Kurdyka-{\L}ojasiewicz extension for functions definable in
an o-minimal structure~\cite{bolte}; this is the form of the hypothesis in which Absil, Mahony and
Andrews prove convergence of a descent method with an Armijo line search to a single point~\cite{absil}.
Write $S_j=F(z^{(j)})-F^\star$ and $\Delta_j=\|z^{(j+1)}-z^{(j)}\|$. Theorem~\ref{thm:descent} with the
step floor and the inequality above give
\begin{equation}
S_j-S_{j+1}\;\ge\;10^{-4}\alpha_{\min}\,\rho(z^{(j)})^2\;\ge\;10^{-4}\alpha_{\min}c_0^2\,S_j^{2\vartheta},
\label{eq:decrease-kl}
\end{equation}
so $S_j$ decreases to $0$, and $S_j=O\big(j^{-1/(2\vartheta-1)}\big)$, geometrically when
$\vartheta=\tfrac12$. The left inequality of \eqref{eq:decrease-kl} bounds the residual by the
decrease, $\rho(z^{(j)})\le10^{2}\alpha_{\min}^{-1/2}(S_j-S_{j+1})^{1/2}$, and the step ceiling bounds
the displacement by the residual, $\Delta_j\le\alpha_{\max}\sqrt n\,\rho(z^{(j)})$, so
\begin{equation}
\Delta_j\;\le\;K\,(S_j-S_{j+1})^{1/2},\qquad K=10^{2}\alpha_{\max}\sqrt n\,\alpha_{\min}^{-1/2}.
\label{eq:tv-bound}
\end{equation}
The right inequality of \eqref{eq:decrease-kl} and the mean value theorem give
$S_j^{1-\vartheta}-S_{j+1}^{1-\vartheta}\ge(1-\vartheta)S_j^{-\vartheta}(S_j-S_{j+1})
\ge(1-\vartheta)10^{-4}\alpha_{\min}c_0^2S_j^{\vartheta}$, and combining this with
\eqref{eq:decrease-kl} to bound $S_j^{\vartheta}(S_j-S_{j+1})^{-1/2}$ yields
$(S_j-S_{j+1})^{1/2}\le K'(S_j^{1-\vartheta}-S_{j+1}^{1-\vartheta})$ for a constant $K'$ depending only
on $\vartheta$, $\alpha_{\min}$ and $c_0$. The sum over $j$ therefore telescopes,
$\sum_j\Delta_j\le KK'\,S_0^{1-\vartheta}<\infty$, so the iterates are Cauchy. Their limit $z^\star$
satisfies $\rho(z^\star)=0$ because $\rho(z^{(j)})\to0$ and $\rho$ is continuous, and
$\|z^{(j)}-z^\star\|\le\sum_{k\ge j}\Delta_k=O(S_j^{1-\vartheta})=O\big(j^{-(1-\vartheta)/(2\vartheta-1)}\big)$
gives the stated rate.
\end{proof}

\begin{remark}[What the theorem does not say]
Convergence to \emph{a} stationary point is not convergence to the global minimum, and a nonconvex
energy may have several stationary points with different basins. Different initial states, and
different sensory boundary data, may settle differently. The system is a family of equilibria indexed
by inputs, clamps and parameters, which is the intended semantics, and Theorem~\ref{thm:convergence} is
what licenses reading the settled state as \emph{the} answer to the current query.
\end{remark}

\section{Inference and learning are the same operator}
\label{sec:oneop}

\subsection{The anchored objective}
A learning solve fixes supplied output targets $y$ as clamps and makes the parameters eligible, adding
a proximal anchor at the parameters held immediately before the experience:
\begin{equation}
F_{\text{learn}}(x,\theta)=\tfrac12\sum_i\big(x_i-p_i(\theta)\big)^2
+\tfrac{\lambda}{2}\|x\|^2
+\tfrac{\rho}{2}\,\|\theta-\theta_a\|^2 .
\label{eq:anchored}
\end{equation}
The anchor $\theta_a$ is fixed for the whole solve and charged once. In a batch of $B$ labelled
experiences, each row carries private state and the rows share one parameter vector and one anchor,
\begin{equation}
F_{\text{batch}}=\frac1B\sum_{b=1}^{B}E_b(x_b,\theta)+\frac{\rho}{2}\|\theta-\theta_a\|^2 ,
\label{eq:batch}
\end{equation}
and qualification applies the projected stationarity test to each row's unaveraged state gradient and
to the shared parameter gradient of the complete mean-plus-anchor objective, so that the tolerance does
not depend on $B$.

\subsection{The role of the anchor}
\begin{theorem}[Well-posedness of the parameter block]
\label{thm:anchor}
Fix a bounded box for $\theta$ and let $M=\sup_{\theta}\|\nabla_\theta^2 E\|$ where $E$ is the data
part of \eqref{eq:anchored} on that box. Then
$\nabla_\theta^2 F_{\text{learn}}\succeq(\rho-M)I$. In particular, if $\rho>M$ the parameter block is
$(\rho-M)$-strongly convex, and projected gradient descent on it with step $1/L_\theta$ contracts with
factor $1-(\rho-M)/L_\theta<1$, where $L_\theta$ is the Lipschitz constant of the parameter gradient on
the box.
\end{theorem}
\begin{proof}
The Hessian of the anchor term is $\rho I$, and the Hessian of the data term is bounded below by
$-MI$ on the box by definition of $M$. Summing gives the first claim. Strong convexity with modulus
$\rho-M$ and $L_\theta$-Lipschitz gradient gives the standard contraction rate
$1-(\rho-M)/L_\theta$ for projected gradient descent on the parameter block, and choosing $\rho>M$
makes the rate strictly less than one.
\end{proof}

\begin{remark}
Equation~\eqref{eq:anchored} is a maximum-a-posteriori objective with a Gaussian prior of precision
$\rho$ at the pre-experience parameters. Theorem~\ref{thm:anchor} makes the prior's role exact: it is
what bounds the influence of any single experience on the parameters, which is the property required
for continual learning and is absent from an unanchored least-squares fit. The corresponding design
rule sets the anchor weight above the local curvature of the data term.
\end{remark}

\subsection{One operator, two blocks}
\begin{theorem}[No train/inference distinction]
\label{thm:oneop}
Let $\cX$ be the free state coordinates and $\Theta$ the parameter coordinates. Define
$G(x,\theta)=\nabla_{(x,\theta)}F_{\text{learn}}$ and let $G_x$ be its state block. Then (i) $G_x$ does
not depend on whether $\Theta$ is eligible, and equals the state gradient of the query objective
$F_{\text{query}}=E$; (ii) the repair operator \eqref{eq:repair} applied to $\cX$ alone is the query
solve, and applied to $\cX\cup\Theta$ is the learning solve; (iii) as $\rho\to\infty$ the learning
solve converges to the query solve with the parameter block frozen.
\end{theorem}
\begin{proof}
(i) The anchor term $\tfrac{\rho}{2}\|\theta-\theta_a\|^2$ has zero derivative with respect to $x$
and adds only $\rho(\theta-\theta_a)$ to the parameter gradient, so the state block of
$\nabla F_{\text{learn}}$ equals $\nabla_x E$ identically, and $F_{\text{learn}}-E$ is a function of
$\theta$ alone. (ii) Eligibility enters \eqref{eq:repair}
only through which coordinates are projected and tested in \eqref{eq:residual}, and the Armijo test
compares two values of the same objective, so an additive constant in the objective changes no
acceptance decision: restricted to $\cX$, the two solves produce the same iterates pointwise, and the
two solves are one map with two coordinate index sets. (iii) At a stationary point of the learning solve
the parameter block satisfies $\theta-\theta_a=-\rho^{-1}\nabla_\theta E(\theta,x)$ exactly, so
$\|\theta-\theta_a\|_\infty\le\rho^{-1}\sup_{\text{box}}\|\nabla_\theta E\|_\infty=O(1/\rho)$, the
supremum being finite because a $C^1$ function has bounded gradient on a compact box. The parameter
block therefore collapses onto the anchor, and the state subproblem, whose solution is unique and
continuous in $\theta$ on the nondegenerate branch, converges to the frozen-parameter problem at
$\theta_a$.
\end{proof}

Theorem~\ref{thm:oneop} formalises the statement that the system has no training phase: there is one
settling operator, one energy and one qualification test, and the distinction between querying the
system and teaching it consists solely in which coordinates the operator may move.

\subsection{Equilibrium propagation}
The learning contrast of a patch net is equilibrium propagation~\cite{scellier2017}. We restate the
result and prove it for the energy \eqref{eq:energy}, because the biological and evolutionary
arguments in \S\S\ref{sec:evo}--\ref{sec:bio} depend on its exact statement. The source states the
lemma for a fixed point that exists and is nondegenerate, and it treats the free fixed point as a
selected branch of the set of local minima. In the patch net the existence and the branch are
consequences of the settling rule rather than assumptions: the sublevel sets are bounded by
Theorem~\ref{thm:coercive} and Theorem~\ref{thm:convergence} gives convergence to a single stationary
point from a given initial state. Positive definiteness of the Hessian at that point remains a
hypothesis. Recent treatments make the same assumption explicitly, list convergence to an equilibrium
as a limitation of the method rather than a property of it~\cite{laborieux2024}, and take the loss to
be independent of the parameters~\cite{litman2025}, which is the convention used here.

Let $C(x,y)$ be a differentiable cost, let
\begin{equation}
F(x,\beta)=E(x)+\beta\,C(x,y)
\label{eq:total}
\end{equation}
be the total energy, let $x^0=\arg\min_x F(x,0)$ be the free state and
$x^\beta=\arg\min_x F(x,\beta)$ the nudged state. Define the one-sided contrast
\begin{equation}
D(\theta,\beta)=\frac1\beta\Big[\partial_\theta E(\theta,x^\beta)-\partial_\theta E(\theta,x^0)\Big].
\label{eq:contrast}
\end{equation}

\begin{theorem}[Equilibrium propagation, fundamental lemma]
\label{thm:ep}
Assume $E$ and $C$ are $C^2$, that the cost is independent of the parameters,
$\partial_\theta C=0$, as it is for the squared output cost of a patch net, that $x^0$ is a
nondegenerate local minimum of $E$ so that
$H=\nabla_x^2E(\theta,x^0)$ is positive definite, and that $x^\beta$ depends $C^1$ on $(\theta,\beta)$
in a neighbourhood of $(0,\theta_0)$, which the implicit function theorem supplies. Then
\begin{equation}
\lim_{\beta\to0^+}D(\theta,\beta)=\frac{\partial L}{\partial\theta}(\theta),
\qquad L(\theta)=C\big(x^0(\theta),y\big),
\end{equation}
with residual $O(\beta)$. Consequently the update $\Delta\theta=-\eta\,D(\theta,\beta)$ is one step of
stochastic gradient descent on $L$ with learning rate $\eta$, up to $O(\beta)$.
\end{theorem}
\begin{proof}
Write $g(\theta,\beta,x)=\nabla_xE(\theta,x)+\beta\nabla_xC(x,y)$, so that the two phases satisfy
$g(\theta,0,x^0)=0$ and $g(\theta,\beta,x^\beta)=0$.
Differentiating the second identity in $\beta$ at $\beta=0$,
\[
H\,\partial_\beta x^\beta\big|_{\beta=0}+\nabla_xC(x^0,y)=0
\quad\Longrightarrow\quad
\partial_\beta x^\beta\big|_{\beta=0}=-H^{-1}\nabla_xC(x^0,y).
\]
Differentiating the first identity in $\theta$ gives
$\partial_\theta x^0=-H^{-1}\partial^2_{x\theta}E$. By the chain rule,
\[
\frac{\partial L}{\partial\theta}
=\big(\partial_\theta x^0\big)^{\!\top}\nabla_xC(x^0,y)
=-\big(\partial^2_{x\theta}E\big)^{\!\top}H^{-1}\nabla_xC(x^0,y),
\]
where $H$ is symmetric because $E\in C^2$. On the other side, the implicit function theorem makes
$\beta\mapsto\partial_\theta E(\theta,x^\beta)$ differentiable at $\beta=0$ with derivative
$\partial^2_{\theta x}E\cdot\partial_\beta x^\beta|_{\beta=0}
=-\partial^2_{\theta x}E\,H^{-1}\nabla_xC(x^0,y)$. The two expressions agree because
$\partial^2_{\theta x}E$ and $\partial^2_{x\theta}E$ are transposes and $H^{-1}$ is symmetric. The
residual is the first-order Taylor remainder of that map, which is $O(\beta)$ under the stated
regularity.
\end{proof}

\begin{corollary}[Locality]
The contrast for coordinate $\theta_k$ uses only the two values of $\partial_{\theta_k}E$ at the two
states, both of which are available at the location of that parameter. No transposed weight matrix, no
stored forward pass and no separate error circuit is required.
\end{corollary}

\begin{remark}[The general cost]
If the cost does depend on the parameters, as it does when a regulariser is written into $C$, the
objective acquires the extra term $\partial_\theta C(\theta,x^0)$ and the contrast alone is no longer
the gradient. This is the general statement of the source, whose equivalent form keeps that
term~\cite{scellier2017}; the patch net excludes it by construction.
\end{remark}

\begin{corollary}[Equality with backpropagation through time]
For discrete-time dynamics whose transition function derives from a primitive, the first $K$ state
updates of the second phase coincide, up to sign, with the first $K$ error terms that backpropagation
through time computes on the same graph, and in
continuous time the second-phase temporal derivative, rescaled by $1/\beta$, reproduces the error
derivatives computed iteratively by recurrent backpropagation~\cite{scellier2019,ernoult2019}. Stopping
the second phase early is truncated recurrent backpropagation.
\end{corollary}

\subsection{The centred contrast is second order}
\label{sec:centred}
The contrast \eqref{eq:contrast} is a finite difference in $\beta$. The source evaluates the one-sided
form, taking the limit $\beta\to0^{+}$ and using a random sign of $\beta$ in the second
phase~\cite{scellier2017}. A one-sided difference is first-order
accurate; a centred difference, which contrasts a $+\beta$ and a $-\beta$ nudge from the same free
state, is second-order. Implementations use the centred form.

\begin{theorem}[Order of the centred contrast]
\label{thm:centred}
Let $f(\beta)=\partial_\theta E(\theta,x^\beta)$ and suppose $f$ is $C^3$ near $\beta=0$. Define the
centred contrast
\begin{equation}
D_{\text{ctr}}(\theta,\beta)=\frac{1}{2\beta}\Big[\partial_\theta E(\theta,x^{+\beta})-\partial_\theta E(\theta,x^{-\beta})\Big].
\label{eq:centred}
\end{equation}
Then $D_{\text{ctr}}=\partial_\theta L+O(\beta^2)$, whereas $D=\partial_\theta L+O(\beta)$.
\end{theorem}
\begin{proof}
By Theorem~\ref{thm:ep}, $f$ is the restriction to the equilibrium branch of a $C^2$ map, and $x^\beta$
is the minimiser of $E+\beta C$ for both signs of $\beta$ by the same implicit equation with $\beta$
replaced by $-\beta$. The centred difference $(f(\beta)-f(-\beta))/(2\beta)$ equals $f'(0)$ plus the
third-order Taylor remainder of $f$, which is $O(\beta^2)$ because the even-order terms cancel by
antisymmetry of the difference. The one-sided difference $(f(\beta)-f(0))/\beta$ equals $f'(0)$ plus
the second-order remainder, which is $O(\beta)$. Theorem~\ref{thm:ep} identifies $f'(0)$ with
$\partial_\theta L$.
\end{proof}

The centred form requires one additional settling phase and no additional state storage, since both
nudged phases start from the same free state. The difference in accuracy is substantial:
Example~\ref{ex:chain} reports a factor of $15.5$ at $\beta=4\times10^{-2}$, with fitted orders $1.00$
and $1.94$.

The finite-nudge bias is a known feature of the deterministic rule, and there are two exits from it
that do not go through the centred difference. In a thermal ensemble the contrast of expectations is
the exact gradient of a free-energy difference at every finite $\beta$ and every temperature, so the
bias is absent once the two states are Gibbs samples rather than minimisers~\cite{litman2025}; for
complex-differentiable dynamics a Cauchy integral over an oscillating nudge recovers the exact state
derivative at finite amplitude, which removes the finite-nudge bias and leaves the bias from
asymmetry of the Jacobian to be handled separately~\cite{laborieux2022,laborieux2024}. The centred
difference proved above is the deterministic real-valued variant: it reduces the bias from $O(\beta)$
to $O(\beta^2)$ and does not remove it. It is the form the reference implementation uses, and neither
exact route is available to it without a stochastic ensemble or complex-valued dynamics.

\subsection{Three-factor plasticity realises the contrast}
The contrast is a product of two local activities and one broadcast scalar, which is the structure of a
three-factor plasticity rule: a synapse-local eligibility trace set by pre- and post-synaptic
co-activation, and a third factor that gates the weight change. The reference implementation has this
structure: every synapse records the activity of its own two endpoints in the two phases, a critic
population predicts return from the settled state, and a single neuromodulatory signal broadcasts the
temporal-difference error of that prediction, with the weight change proportional to the product of the
three~\cite{gerstner2018,fremaux2015}. The time constant of the eligibility trace must match the
interval between an action and its outcome, which in the biological literature is on the order of a
second. In the notation of this paper the gate is the nudge and the trace is the pair of settled
phases. The third broadcast factor permits the contrast to be read without a backward pass. A
continuous-time reading, in which the weights move during the relaxation rather than between two
settled states, requires plasticity to be slow compared with settling, which is Criterion 5 of
\S\ref{sec:bio}: otherwise the state being compared has already moved and the contrast in
\eqref{eq:contrast} is not the quantity the lemma bounds.

\section{The evolutionary layer}
\label{sec:evo}

\subsection{Meta-parameters and the lifetime map}
Evolution does not act on the weights of a settling brain within a lifetime, and it need not. Following
the reference design~\cite{cadence,evolution}, the inherited object is a \emph{genome} that determines
the brain's regions, their sizes and their projections. We take
\begin{equation}
\varphi=\big(\text{wiring mask }A_\varphi,\ \text{energy form }E_\varphi,\ \text{plasticity constants }(\eta,\beta,\rho)\big)\in\Phi
\end{equation}
and we require that $\varphi$ \emph{does not contain} the weights learned within a lifetime. A lifetime is
defined by a task stream $\omega$ drawn from $(\Omega,\nu)$ and the EP map
\begin{equation}
\Lambda:\Phi\times\Omega\to\Theta,\qquad
\Lambda(\varphi,\omega)=\theta_K,\qquad
\theta_{k+1}=\theta_k-\eta\,D(\theta_k,\beta),
\label{eq:lifetime}
\end{equation}
with $D$ the contrast of \eqref{eq:contrast}. The effective fitness of a meta-parameter is the
\emph{post-learning} fitness
\begin{equation}
\hat J(\varphi)=\E_{\omega\sim\nu}\big[J\big(\Lambda(\varphi,\omega),\omega\big)\big].
\label{eq:baldwin}
\end{equation}
This is the Baldwinian fitness: the fitness attained after a lifetime of learning, rather than the innate
fitness at birth. A population has access only to realised lifetimes: generation $g$ draws
$\omega_g\sim\nu$, and the field it selects on is the realised fitness
$J_{\omega}(\varphi)=J(\Lambda(\varphi,\omega),\omega)$, not its expectation. The realised slow flow is
therefore
\begin{equation}
\partial_t\rho=\varepsilon\Big[\kappa\,\Delta_\varphi\rho-\nabla_\varphi\!\cdot\big(\rho\,\nabla_\varphi J_{\omega}\big)\Big],
\qquad \omega\sim\nu \text{ resampled at each generation},
\label{eq:realized}
\end{equation}
with $\varepsilon=\tau_{\text{life}}/\tau_{\text{evolution}}\ll1$. Replacing the realised field by its
average gives the Baldwinian flow
\begin{equation}
\partial_t\rho=\varepsilon\Big[\kappa\,\Delta_\varphi\rho-\nabla_\varphi\!\cdot\big(\rho\,\nabla_\varphi\hat J\big)\Big],
\label{eq:replicator}
\end{equation}
and the next subsection proves that \eqref{eq:realized} converges to \eqref{eq:replicator} as
$\varepsilon\to0$. Both are replicator-mutator flows with the drift directed up the fitness gradient, so
their stationary measures concentrate on high fitness.

\subsection{Two-timescale averaging}
\label{sec:averaging}
\begin{theorem}[Averaging to the Baldwinian flow]
\label{thm:averaging}
Assume
\emph{(H1)} for every $\varphi$ in a compact set the fast system \eqref{eq:lifetime} has an attractor
$\mathcal{A}(\varphi)$ attracting at a rate $\lambda_0>0$ uniformly in $\varphi$, with a uniformly
absorbing neighbourhood of the initial states;
\emph{(H2)} the relevant maps are Lipschitz, $J$ is bounded, $\Lambda$ is measurable, and
$\hat J\in C^1(\Phi)$;
\emph{(H3)} the mutation kernel is conservative with a spectral gap;
\emph{(H4)} the task draws are independent with law $\nu$ and the realised gradients
$\nabla_\varphi J_{\omega}$ obey the uniform law of large numbers on compact sets.
Then the slow trajectories of \eqref{eq:realized} converge uniformly on finite time intervals to the
solution of \eqref{eq:replicator} as $\varepsilon\to0$, and the empirical measure of a population of $N$
coupled individuals converges weakly to the same limit as $N\to\infty$ and then $\varepsilon\to0$.
\end{theorem}
\begin{proof}[Proof sketch]
Write the pair as a singularly perturbed system $\dot\varphi=\varepsilon f(\varphi,\theta,\omega)$,
$\dot\theta=g(\varphi,\theta)+\sigma\xi$, with $f$ the realised fitness gradient. Assumption (H1) makes
the fast flow converge to a $\varphi$-indexed invariant manifold, and Tikhonov's theorem for uniformly
asymptotically stable fast dynamics gives uniform convergence on $[0,T]$ to the restricted equation,
whose drift is the ergodic average of $f$ along the fast flow. By (H4) that average is
$\E_{\omega}\nabla_\varphi J_{\omega}$ uniformly up to $o(1)$, and by \eqref{eq:baldwin}, (H2) and
differentiation under the expectation it equals $\nabla_\varphi\hat J$. Replacing the realised field by
its average is the stochastic-averaging step, and the deterministic part of its error is
$O(\varepsilon)$; the residual random fluctuation of \eqref{eq:realized} about \eqref{eq:replicator} is
of larger order and is not controlled by this rate, which is why the finite-population spread in
\S\ref{ex:evo} exceeds the mean-field deviation. The measure-valued
statement follows by compactness of the empirical measures under the standard propagation-of-chaos
argument, with (H2) giving continuity of the drift functional in the weak topology and (H3) preventing the
diffusion from degenerating.
\end{proof}

\begin{remark}
Theorem~\ref{thm:averaging} gives the Baldwin effect a precise form. Whether evolution acts on the
innate fitness or on the post-learning fitness is not a modelling choice; it is determined by the
time-scale ratio, and the theorem establishes that the limit is the post-learning fitness. The
deterministic error of the approximation is $O(\varepsilon)$, which quantifies
how rapid a lifetime must be relative to a generation and is the quantity the averaging-order prediction
of \S\ref{sec:limits} measures.
\end{remark}

\subsection{Decomposition}
\begin{proposition}[Price-type decomposition]
\label{prop:price}
Write $\hat J=\hat J_0+\Delta\hat J_{\text{learn}}$, where $\hat J_0$ is the fitness with learning
frozen, so that $\Delta\hat J_{\text{learn}}=\hat J-\hat J_0$ is what a lifetime of learning adds. Let
$w$ be the relative fitness of the transmission step, normalised so that $\bar w=1$. Then the
per-generation change of mean fitness decomposes as
\begin{equation}
\overline{\Delta\hat J}=
\underbrace{\Cov_\rho(w,\hat J_0)}_{\text{innate selection}}
+\underbrace{\Cov_\rho(w,\Delta\hat J_{\text{learn}})}_{\text{plasticity selection}}
+\underbrace{\E_\rho\big[w\,\Delta\hat J_{\text{learn}}\big]}_{\text{transmission}}
+\underbrace{\mathcal{M}}_{\text{mutation}} .
\label{eq:price}
\end{equation}
At a stationary population $\rho^\star$ of \eqref{eq:replicator} the gradient condition
$\nabla\hat J=\kappa\nabla\log(\rho^\star/\rho_{\text{ref}})$ holds by
Theorem~\ref{thm:dissipation}, so the mean fitness is stationary and the mutation term balances the
selection and transmission terms rather than being a small correction to them; the stationary measure
is the Gibbs form $\rho^\star\propto\rho_{\text{ref}}e^{\hat J/\kappa}$, which concentrates on high
fitness.
\end{proposition}
\begin{proof}
\begin{sloppypar}
The identity is the standard Price covariance-and-transmission decomposition applied to the trait
$z=\hat J$ and split by linearity into its innate and learned components, with the transmission term
carrying the relative fitness as in the original derivation~\cite{price,frank}; the normalisation
$\bar w=1$ removes the divisor of the usual form. The stationarity claim follows from
Theorem~\ref{thm:dissipation}: at a stationary measure $\delta\cF/\delta\rho$ is constant, which by
\eqref{eq:replicator} is $\kappa\nabla\log(\rho^\star/\rho_{\text{ref}})=\nabla\hat J$, and
$\overline{\Delta\hat J}=0$ because the measure does not change. The two covariance terms do not vanish
individually; it is their sum together with the transmission term that mutation balances, which is the
mutation-selection balance of the flow.
\end{sloppypar}
\end{proof}

\subsection{Population dissipation}
\begin{theorem}[Free-energy dissipation]
\label{thm:dissipation}
Let
\begin{equation}
\cF(\rho)=-\langle\hat J\rangle_\rho+\kappa\!\int_\Phi\rho\log\frac{\rho}{\rho_{\text{ref}}}\,\mathrm d\varphi,
\label{eq:freeenergy}
\end{equation}
the negated mean fitness plus $\kappa$ times the relative entropy to a fixed reference measure
$\rho_{\text{ref}}$. Solutions of \eqref{eq:replicator} satisfy
\begin{equation}
\frac{\mathrm d}{\mathrm dt}\cF(\rho_t)=-\int_\Phi\rho\,\Big|\nabla_\varphi\Big(\frac{\delta\cF}{\delta\rho}\Big)\Big|^2\mathrm d\varphi\le0,
\end{equation}
with equality if and only if $\rho=\rho^\star\propto\rho_{\text{ref}}e^{\hat J/\kappa}$. The four layers
of the system carry a nested Lyapunov family $\{E,\,L,\,\cF_{\text{projection}},\,\cF\}$, one per time
scale.
\end{theorem}
\begin{proof}
Equation~\eqref{eq:replicator} is the Wasserstein-2 gradient flow of $\cF$, so
$\partial_t\rho=\nabla\cdot(\rho\nabla\frac{\delta\cF}{\delta\rho})$. Substituting into
$\frac{\mathrm d}{\mathrm dt}\cF=\int\frac{\delta\cF}{\delta\rho}\partial_t\rho$ and integrating by
parts gives the identity; the additive constant in $\delta\cF/\delta\rho$ contributes nothing because
$\rho$ is a probability measure and $\int\partial_t\rho=0$. The equilibrium condition
$-\nabla\hat J+\kappa\nabla\log(\rho/\rho_{\text{ref}})=0$ gives
$\rho^\star\propto\rho_{\text{ref}}e^{\hat J/\kappa}$. The nested family
is Theorem~\ref{thm:descent}, \ref{thm:ep} and \ref{thm:dissipation} together with the
information-projection inequality at the consolidation scale.
\end{proof}

\begin{remark}[Energy accounting]
Theorem~\ref{thm:dissipation} is a statement about structure rather than about energy in joules. The
thermodynamic account is separate, and it splits in the same direction as the structural one: for an
equilibrium physical network the free energy is independent of the input and of the parameters, so
quasi-static inference requires no work at all, while writing the parameters carries an irreducible
cost of a few $k_BT$ each, which survives the quasi-static limit and is set by the memory rather than by
the arithmetic~\cite{thermo_bounds}; an analog substrate is the natural place for the free part, since
it relaxes by its own dynamics~\cite{kendall2020}. The consequence for this architecture is
structural, not a measured energy saving: a claim that it reduces the energy cost of
learning would require separate measurement.
\end{remark}

\section{Biological compliance}
\label{sec:bio}

\subsection{Six formal criteria}
A construction can be described as biologically motivated without being checked against biology. To
make such a check possible, we write the basic regularities of biological evolution as formal
conditions on the permitted dynamics.
\begin{criterion}[Blind variation]
The mutation operator $M(\varphi'\mid\varphi)$ is statistically independent of the fitness $\hat J$ and
of the task distribution $\nu$.
\end{criterion}
\begin{criterion}[Causal and post hoc selection]
Fitness is measured only after a phenotype has performed, and selection acts only through reweighting
of the population measure; it never modifies an individual's internal state during its lifetime.
\end{criterion}
\begin{criterion}[Heredity and the Weismann barrier]
The heritable object is bounded in fidelity, and changes acquired within a lifetime do not enter the
germ line by default.
\end{criterion}
\begin{criterion}[Thermodynamic constraint]
Every layer dissipates non-negative free energy, and no mechanism obtains fitness at zero cost.
\end{criterion}
\begin{criterion}[Time-scale separation]
Settling is fast, lifetime learning is intermediate and evolution is slow, so that learning does not
constitute a second, implicit selection channel.
\end{criterion}
\begin{criterion}[Redundancy]
Phenotype is robust to parameter perturbation, with degenerate pathways supplying damage tolerance.
\end{criterion}

\begin{theorem}[Compliance]
\label{thm:compliance}
Let the system be that of \S\S\ref{sec:settling}--\ref{sec:evo} with heritable object $\varphi$ and
lifetime map \eqref{eq:lifetime}. Then (i) if the mutation operator is independent of $\hat J$ and
$\nu$, selection acts through $\rho$ alone after the life, and only $\varphi$ is inherited, the system
satisfies Criteria 1 to 6; (ii) if the inheritance restriction is dropped so that $\theta_K$ is
transmitted, the system violates Criterion 3, and there exist task distributions $\nu$ on which the
unrestricted system has strictly lower stationary fitness.
\end{theorem}
\begin{proof}
(i) Independence of the mutation operator is Criterion 1 by construction. Selection acts only through
the population measure in \eqref{eq:replicator} and the fitness is the post-life quantity
\eqref{eq:baldwin}, which is Criterion 2. Inheriting $\varphi$ and not $\theta$ is Criterion 3.
Theorem~\ref{thm:dissipation} gives non-negative dissipation with a strict decrease away from
equilibrium, and the thermodynamic bound of the remark above supplies the floor for the learning
layer, which is Criterion 4; the variational reading of free energy as an organising principle for
biological systems is the free-energy principle~\cite{friston}. The scale separation $\tau_{\text{settle}}\ll\tau_{\text{life}}\ll
\tau_{\text{evolution}}$ is the hypothesis of Theorem~\ref{thm:averaging}, which is Criterion 5.
Sparse wiring with high clustering and short paths, the structure measured in the complete
nematode connectome~\cite{connectome}, is the substrate in which the degenerate pathways of
Criterion 6 are available, so the architecture does not preclude them; a measured statement of
robustness under parameter perturbation is not claimed here, and \S\ref{sec:limits} records it as an
open problem.
(ii) Transmitting $\theta_K$ makes an acquired change heritable, which is the negation of Criterion 3.
For the adaptive claim, take the two-point task distribution $\nu=\tfrac12\delta_{-1}+\tfrac12\delta_{+1}$
with $\omega$ alternating across generations, a single linear patch whose clamped target is $\omega$, and
the quadratic cost $C=\tfrac12(\theta-\omega)^2$, so that a lifetime long enough to reach the stationary
point of \eqref{eq:anchored} is the anchored least-squares solution
$\Lambda(\varphi,\omega)=(\omega+\rho\theta_a)/(1+\rho)$ with fitness $J=-\tfrac12(\Lambda-\omega)^2$. The
restricted system starts every lifetime from the $\varphi$-determined initial parameter, which by the
symmetry of $\nu$ may be taken to be $0$, so $\Lambda=\omega/(1+\rho)$ and its stationary fitness is
$\hat J_{\text{rest}}=-\tfrac12\big(\rho/(1+\rho)\big)^2$, the same for both tasks. The unrestricted
system inherits $\theta_K$ of the previous generation, so its anchor is mismatched to the new task and
the inherited value alternates between $\pm c$ with $c=1/(1+2\rho)$ at the steady state; a lifetime on a
task $\omega$ then has $\Lambda-\omega=-\rho(1+c)/(1+\rho)$ and fitness
$\hat J_{\text{unrest}}=-\tfrac12\big(\rho(1+c)/(1+\rho)\big)^2=(1+c)^2\,\hat J_{\text{rest}}<
\hat J_{\text{rest}}$, because $c>0$ whenever $\rho>0$. The two regimes exchange rank when the task
distribution is narrow and stable, where the inherited weights are nearly correct and no re-learning is
needed, which is the ordering reported for Baldwinian against Lamarckian
meta-learning~\cite{baldwin}.
\end{proof}

\begin{remark}
Part (ii) does not assert that acquired inheritance is always disadvantageous. It asserts that the two
regimes are separated by a boundary determined by the rate at which the task distribution changes, and
that the biological default lies on the side of the boundary at which not inheriting is preferable.
The construction is of interest because it parameterises the choice rather than prescribing it.
\end{remark}

\subsection{Limits of the correspondence}
Three limitations are stated explicitly. The units are rate units with a smooth prediction function
rather than spiking neurons, as the reference library states of itself. The body model is simplified.
The averaging theorem is a statement in the limit of a large population and a large time-scale ratio,
whereas a biological lineage is finite and subject to drift.

\section{Worked examples}
\label{sec:examples}

All numbers in this section are produced by the two scripts listed in Appendix~\ref{app:repro}. The
first implements the patch law \eqref{eq:patch} directly; the second imports the reference library and
compares.

\subsection{An exactly solvable instantiation}
\label{ex:chain}
Take one patch reading one sensory value $u$ through one weight $w$ with bias $b$, with the activity
prior switched off, and take the squared cost $C(x,y)=\tfrac12(x-y)^2$. Then
$t=\tanh(b+wu)$ and $E=\tfrac12(x-t)^2$. The nudged state is obtained by solving a quadratic,
$x^\beta=(t+\beta y)/(1+\beta)$, and $x^{-\beta}=(t-\beta y)/(1-\beta)$. The parameter gradient is
$\partial_wE\big|_x=-(x-t)\,(1-t^2)u$, so
\[
f(\beta)-f(0)=-\frac{\beta(y-t)(1-t^2)u}{1+\beta},
\qquad
f(\beta)-f(-\beta)=-\frac{2\beta(y-t)(1-t^2)u}{1-\beta^2}.
\]
Since $\partial L/\partial w=(t-y)(1-t^2)u$,
\begin{equation}
\boxed{\;D_{\text{one}}=\frac{\partial L}{\partial w}\cdot\frac{1}{1+\beta},
\qquad
D_{\text{ctr}}=\frac{\partial L}{\partial w}\cdot\frac{1}{1-\beta^{2}}.\;}
\label{eq:closedform}
\end{equation}
The one-sided contrast is the gradient scaled by $(1+\beta)^{-1}$ and the centred contrast by
$(1-\beta^2)^{-1}$, so both are exact in closed form and the centred error is the square of the
one-sided error to leading order. With $w=0.9$, $b=0.2$, $u=0.55$ and $y=0.7$ we have
$t=0.6011844863$ and $\partial L/\partial w=-0.0347057345$. Table~\ref{tab:closedform} compares the
numerically computed contrasts against \eqref{eq:closedform}; the largest deviation is
$4.7\times10^{-11}$.

\begin{table}[t]
\centering
\small
\begin{tabular}{rrrr}
\toprule
$\beta$ & one-sided (numeric) & one-sided (closed form) & centred (numeric)\\
\midrule
$0.100$ & $-0.03155067$ & $-0.03155067$ & $-0.03505630$\\
$0.050$ & $-0.03305308$ & $-0.03305308$ & $-0.03479272$\\
$0.020$ & $-0.03402523$ & $-0.03402523$ & $-0.03471962$\\
$0.010$ & $-0.03436211$ & $-0.03436211$ & $-0.03470921$\\
\bottomrule
\end{tabular}
\caption{The exactly solvable case of \S\ref{ex:chain}. The centred column matches
$\partial L/\partial w\cdot(1-\beta^2)^{-1}$ to the same $4.7\times10^{-11}$; the target value is
$\partial L/\partial w=-0.0347057345$.}
\label{tab:closedform}
\end{table}

\subsection{Settling, coercion and the role of the prior}
The second example is a three-patch graph, $p_0$ reading two sensory values, $p_1$ reading $x_0$, a
sensory value and $x_2$ (a feedback edge), and $p_2$ reading $x_1$, with $\lambda=0.01$. The repair
trajectory decreases the energy at every accepted sweep over $2001$ sweeps, from $0.22303138$ to
$0.00396407$, and terminates with a projected stationarity residual of $4.297\times10^{-13}$.

The coercivity bound of Theorem~\ref{thm:coercive} is checked directly: at a random state with
$\|x\|_\infty=15.98$, $E=209.1931$ while $\tfrac12\sum_i((|x_i|-1)^+)^2=205.0945$, so the bound holds
with slack; the same lower bound applies at $\lambda=0$ and at $\lambda=0.01$, which is the sense in
which it is independent of the activity prior. The role of the prior appears in the solution obtained.
With
$\lambda=0$ the solve reaches the consistent configuration exactly: $E^\star=0$, largest prediction
error $3.3\times10^{-16}$. With $\lambda=0.01$ it reaches a compromise: $E^\star=0.00396407$ with
largest prediction error $1.4\times10^{-2}$. Both are stationary points corresponding to different
objectives. This is the quantitative form of the statement that qualification measures stationarity and
not prediction error.

\subsection{Order of the contrast}
On the same three-patch graph with the feedback edge, the contrast is compared against the true
gradient of the free-phase loss, obtained by central differences of the settled free state, which is
the implicit-function-theorem gradient and not an approximation to it. Table~\ref{tab:order} reports
both contrasts. Halving $\beta$ halves the one-sided error and quarters the centred error; the fitted
orders are $1.00$ and $1.94$. At $\beta=4\times10^{-2}$ the centred contrast is more accurate by a
factor of $15.5$.

\begin{table}[t]
\centering
\small
\begin{tabular}{rcc}
\toprule
$\beta$ & one-sided $\|D-\partial L/\partial\theta\|$ & centred $\|D-\partial L/\partial\theta\|$\\
\midrule
$4.00\times10^{-2}$ & $7.968\times10^{-2}$ & $5.125\times10^{-3}$\\
$2.00\times10^{-2}$ & $4.096\times10^{-2}$ & $1.283\times10^{-3}$\\
$1.00\times10^{-2}$ & $2.076\times10^{-2}$ & $3.214\times10^{-4}$\\
$5.00\times10^{-3}$ & $1.045\times10^{-2}$ & $8.111\times10^{-5}$\\
$2.50\times10^{-3}$ & $5.245\times10^{-3}$ & $2.116\times10^{-5}$\\
\midrule
fitted order & $1.00$ & $1.94$\\
\bottomrule
\end{tabular}
\caption{Contrast error against the exact gradient of the free-phase loss on a three-patch graph with
a feedback edge. The centred form used by implementations reduces the error by a factor of $248$ at
$\beta=2.5\times10^{-3}$.}
\label{tab:order}
\end{table}

\subsection{One operator and convergent learning}
On the same graph and at the same point $(x,\theta)$, the state block of the gradient is identical
whether or not the parameters are eligible: the measured difference is exactly $0.0$, as
Theorem~\ref{thm:oneop} requires. A learning solve with a clamped output target decreases its anchored
objective monotonically from $0.50822925$ to $0.03354746$ and terminates with a projected residual of
$2.459\times10^{-15}$. Presenting six successive experiences to the same operator reduces the
free-phase cost of the clamped target from $9.036\times10^{-3}$ to $2.649\times10^{-4}$, which is the
patch-net instance of the source's statement that for small $\beta>0$ the nudged fixed point has lower
cost than the free fixed point~\cite{scellier2017}, while the
parameter displacement accumulates steadily. This is characteristic of a system that learns from each
experience without discarding what it has already acquired.

\subsection{Agreement with the reference implementation}
\label{ex:crosscheck}
The second script imports the reference library at version 0.80.0 and compares, on the same graph,
state and parameters: the energy, the prediction, the error, and the state, weight and bias gradients.
The largest absolute difference over all of them is $1.1\times10^{-16}$, which is float64 rounding.
Driving both implementations from the same initial state with the same budget produces terminal states
agreeing to $1.018\times10^{-12}$ and terminal energies agreeing to $2.2\times10^{-19}$. A settlement
through the library's public composition interface qualifies with a stationarity residual of
$2.8\times10^{-7}$ after seven accepted sweeps. The mathematics in this paper therefore describes the
implementation itself rather than an idealisation of it.

\subsection{The evolutionary layer}
\label{ex:evo}
For the evolutionary example we use scalar meta-parameter $\varphi$ standing for plasticity, a lifetime
of three local updates on a task drawn from the environment, and a metabolic price $\lambda_\varphi$
for plasticity. The post-learning fitness is
$\hat J(\varphi)=-\tfrac12(1-\varphi)^6-\lambda_\varphi\varphi^2$ in closed form, whose maximum for
$\lambda_\varphi=0.35$ is at $\varphi^\star=0.3913$; the estimated maximum on a grid of $161$ values
agrees with the closed form to $7.2\times10^{-3}$, which is the Monte-Carlo error of the estimate. The
optimum being interior and positive is the statement that learning is retained; a value of zero would
correspond to a lineage that has eliminated plasticity.

Solving \eqref{eq:replicator} on the grid with mutation strength $\kappa=0.02$ decreases the population
free energy monotonically from $0.379150$ to $0.109856$, and the population measure approaches the
Gibbs form $e^{\hat J/\kappa}$ to a total-variation distance of $0.0040$. An individual-based
simulation with $N=4000$ individuals over $400$ generations, in which each individual performs its own
lifetime of updates and is selected on a noisy estimate of its fitness, reaches mean
$0.4262\pm0.1402$ against the mean-field prediction $0.4001$. The mean-field flow is thus the limit of
the finite population, and the residual spread is the finite-population drift that
Theorem~\ref{thm:averaging} does not claim to capture at first order.

\section{Limits and falsifiable predictions}
\label{sec:limits}

\subsection{Open problems}
\begin{enumerate}
\item \textbf{Uniform hyperbolicity of the fast flow.} Theorem~\ref{thm:averaging} assumes the lifetime
dynamics attract uniformly across meta-parameters. The settling theorem, Theorem~\ref{thm:convergence},
guarantees convergence for each fixed configuration, but not uniformly with a rate bounded away from
zero as the configuration varies. This is the weakest assumption in the argument and the most productive
direction for further work.
\item \textbf{Non-differentiable meta-parameters.} Part of $\varphi$ is combinatorial, namely the wiring
mask. The fitness is then not differentiable in $\varphi$ and the averaging theorem needs a stochastic
approximation or evolution-strategy form, with weaker convergence conditions.
\item \textbf{Estimator noise inside the fitness.} When $\hat J$ is estimated from finitely many tasks,
selection and estimation noise interact across generations, and the mutation-selection balance of
Proposition~\ref{prop:price} is reached only at the nominal stationary point rather than along the
trajectory.
\item \textbf{Consolidation is not yet inside the unified averaging.} Offline reorganisation at an
intermediate time scale is described as a constrained information projection, which is the mechanism
that limits catastrophic forgetting, but its scale separation has not been folded into
Theorem~\ref{thm:averaging}.
\item \textbf{Scale.} Every result in this paper is on graphs of a few patches, and the largest systems
in the reference line of work are small by the standards of contemporary language models. No result
here supports a claim about scale.
\item \textbf{Redundancy is available but unmeasured.} Criterion 6 is met in the sense that sparse
clustered wiring leaves degenerate pathways available, but robustness of a patch net under parameter
perturbation has not been measured on the reference architecture.
\end{enumerate}

\subsection{Predictions}
The following predictions are stated so that they are falsifiable.
\begin{enumerate}
\item \textbf{Order of the centred contrast.} Any patch net whose second phase is an equilibrium branch
of a $C^2$ family should show a centred contrast error scaling as $\beta^2$ and a one-sided error
scaling as $\beta$. A measured first-order centred error refutes Theorem~\ref{thm:centred}.
\item \textbf{Anchor weight above curvature.} For an anchored learning solve, setting $\rho$ above the
local curvature of the data term should remove the oscillatory or divergent parameter trajectories seen
at small $\rho$, and the improvement should track the contraction factor of Theorem~\ref{thm:anchor}.
\item \textbf{Averaging order.} Doubling the lifetime-to-generation ratio should halve the deviation
between the individual-based population and the averaged flow. A measured exponent away from one, for
instance a square-root drift-dominated scaling, refutes the hypotheses of Theorem~\ref{thm:averaging}.
\item \textbf{Weismann boundary.} On a broad or rapidly changing task distribution, the population that
inherits only plasticity should reach strictly higher stationary fitness than the population that also
inherits learned weights, and on a narrow stable distribution the ordering should reverse.
\end{enumerate}

\section{Conclusion}
A learning system requires neither a backward pass nor a phase boundary between training and use. It
requires an energy, a local repair rule that descends that energy, and a single channel through which
the objective enters. We have shown that when the repair is projected gradient descent on the patch
energy used by the reference implementation, the descent is coercive, the iterates converge to a single
stationary point, inference and learning are one operator restricted to different coordinate sets, and
the contrast computed by that operator is second order rather than first. The evolutionary layer is
defined on the same structure and is not merely an analogy: its effective fitness is provably the
post-learning fitness, its compliance with the formal criteria of biological evolution is proved rather
than asserted, and the restriction that acquired weights are not inherited is shown to be advantageous
on the task distributions where the distinction matters. The remaining gaps have been stated, and no
claim about scaling is made.

\appendix

\section{Reproduction}
\label{app:repro}
The evidence in \S\ref{sec:examples} is produced by
\begin{verbatim}
python3 src/verify_eep.py     # Parts A-F, numpy only, Python 3.9+
PYTHONPATH=../cadence/src python3.13 src/cadence_crosscheck.py
\end{verbatim}
The first script asserts each claim and exits non-zero if a claim fails. Its parts are: A, analytic
gradients against central differences ($1.5\times10^{-10}$); B, settling monotonicity, convergence and
coercivity; C, the order of the one-sided and centred contrasts; D, one operator across two coordinate
blocks; E, the evolutionary layer; F, the closed-form case of \S\ref{ex:chain}. The second script
imports the reference library and reproduces \S\ref{ex:crosscheck}. Run logs are archived beside the
scripts.

\section{Structural agreement with observer-patch consensus}
The construction of this paper has the shape of a consensus protocol. A finite set of carriers holds
local state with a declared interface, disagreement is measured by a scalar functional, accepted local
repairs drive that functional strictly down, and the limit is a public normal form that does not depend
on the order of repairs. The physics programme that shares this shape derives its consequences from
three axioms, a finite carrier with a declared boundary, agreement on overlaps, and no structure beyond
the records~\cite{oph}. The correspondence here is structural and not physical. The learning case adds
the second time scale: consensus among patches corresponds to settling and consensus among lineages to
selection, and the two are related by the averaging theorem. This parallel is offered as reason to
expect the structure to be reusable, and not as evidence for either programme.

\begin{thebibliography}{99}
\small
\bibitem{scellier2017} B.~Scellier and Y.~Bengio. Equilibrium propagation: bridging the gap between
energy-based models and backpropagation. \emph{Frontiers in Computational Neuroscience}, 11:24, 2017.
\href{https://arxiv.org/abs/1602.05179}{arXiv:1602.05179}.
\bibitem{scellier2019} B.~Scellier and Y.~Bengio. Equivalence of equilibrium propagation and recurrent
backpropagation. \emph{arXiv:1810.01310}, 2019.
\bibitem{ernoult2019} M.~Ernoult, J.~Grollier, D.~Querlioz, Y.~Bengio and B.~Scellier. Updates of
equilibrium prop match gradients of backprop through time in an RNN with static input.
\emph{NeurIPS}, 2019.
\bibitem{laborieux2022} A.~Laborieux and F.~Zenke. Holomorphic equilibrium propagation computes exact
gradients through finite size oscillations. \emph{NeurIPS}, 2022. \href{https://arxiv.org/abs/2209.00530}{arXiv:2209.00530}.
\bibitem{laborieux2024} A.~Laborieux and F.~Zenke. Improving equilibrium propagation without weight
symmetry through Jacobian homeostasis. \emph{ICLR}, 2024. \href{https://arxiv.org/abs/2309.02214}{arXiv:2309.02214}.
\bibitem{litman2025} E.~Litman. Finite-nudge equilibrium propagation in thermal ensembles.
\href{https://arxiv.org/abs/2511.22024}{arXiv:2511.22024}, 2025.
\bibitem{scellier2018} B.~Scellier, A.~Goyal, J.~Binas, T.~Mesnard and Y.~Bengio. Generalization of
equilibrium propagation to vector field dynamics. \href{https://arxiv.org/abs/1808.04873}{arXiv:1808.04873}, 2018.
\bibitem{kendall2020} J.~Kendall, R.~Pantone, J.~Y.~Zung and Y.~Bengio. Training end-to-end analog
neural networks with equilibrium propagation. \emph{arXiv:2006.01981}, 2020.
\bibitem{cadence} B.~Mueller. Cadence: one continuing equilibrium brain. Software, version 0.80.0,
2026. \url{https://github.com/muellerberndt/cadence}.
\bibitem{manuscript} B.~Mueller. Cadence: a continuing learning system. Manuscript and design
documents, Pragma Research, 2026. \url{https://floatingpragma.io/cadence/}.
\bibitem{evolution} Cadence documentation, \emph{Evolve a brain}: genome, develop, mutate and evolve
across lives. \url{https://github.com/muellerberndt/cadence/blob/main/docs/evolution.md}.
\bibitem{absil} P.-A.~Absil, R.~Mahony and B.~Andrews. Convergence of the iterates of descent methods
for analytic cost functions. \emph{SIAM Journal on Optimization}, 16(2):531--547, 2005.
\bibitem{bolte} J.~Bolte, A.~Daniilidis, A.~Lewis and M.~Shiota. Clarke subgradients of stratifiable
functions. \emph{SIAM Journal on Optimization}, 18(2):556--572, 2007.
\bibitem{lojasiewicz} S.~{\L}ojasiewicz. Sur les trajectoires du gradient d'une fonction analytique.
\emph{Seminari di Geometria 1982--1983}, Universit\`a di Bologna, Istituto di Geometria, Dipartimento
di Matematica, pages 115--117, 1984.
\bibitem{gerstner2018} W.~Gerstner, M.~Lehmann, V.~Liakoni, D.~Corneil and J.~Brea. Eligibility traces
and plasticity on behavioral time scales: experimental support of neoHebbian three-factor learning
rules. \emph{Frontiers in Neural Circuits}, 12:53, 2018.
\bibitem{fremaux2015} N.~Fr\'emaux and W.~Gerstner. Neuromodulated spike-timing-dependent plasticity,
and theory of three-factor learning rules. \emph{Frontiers in Neural Circuits}, 9:85, 2015.
\bibitem{baldwin} C.~Fernando, J.~Syrota and J.~E.~Pompeu. Meta-learning by the Baldwin effect.
\emph{GECCO Companion}, 2018. \href{https://arxiv.org/abs/1806.07917}{arXiv:1806.07917}.
\bibitem{connectome} D.~S.~Bassett and E.~T.~Bullmore. Small-world brain networks revisited.
\emph{The Neuroscientist}, 23(5):499--516, 2017. \href{https://arxiv.org/abs/1608.05665}{arXiv:1608.05665}.
\bibitem{thermo_bounds} A.~V.~Tkachenko. Thermodynamic cost of inference and learning in physical
neural networks. \href{https://arxiv.org/abs/2503.09980}{arXiv:2503.09980}, 2025. Revised 2026.
\bibitem{price} G.~R.~Price. Selection and covariance. \emph{Nature}, 227:520--521, 1970.
\bibitem{frank} S.~A.~Frank. George Price's contributions to evolutionary genetics. \emph{Journal of
Theoretical Biology}, 175(3):373--388, 1995.
\bibitem{hopfield} H.~Ramsauer et al. Hopfield networks is all you need. \emph{ICLR}, 2021.
\bibitem{friston} K.~Friston. A free energy principle for biological systems. \emph{Entropy},
14(11):2100--2121, 2012.
\bibitem{oph} B.~Mueller et al. Observers are all you need; and From observer consensus to standard
physics. Pragma Research, 2026. \url{https://floatingpragma.io/physics}.
\end{thebibliography}

\end{document}
